\documentclass[11pt]{article}
\usepackage[margin=0.9in]{geometry}
\usepackage{amsmath,amssymb,amsthm,booktabs,mathrsfs}
\usepackage[T1]{fontenc}
\usepackage{lmodern,microtype}
\usepackage[colorlinks=true,linkcolor=blue,urlcolor=blue]{hyperref}
\newcommand{\norm}[1]{\left\lVert#1\right\rVert}
\newtheorem{theorem}{Theorem}[section]
\newtheorem{lemma}[theorem]{Lemma}
\newtheorem{proposition}[theorem]{Proposition}
\newtheorem{corollary}[theorem]{Corollary}
\theoremstyle{remark}
\newtheorem{remark}[theorem]{Remark}
\title{Toward a four-dimensional Yang--Mills construction\\
\large Notebook 35: infrared cancellation and physical heat bounds}
\author{Research working document}
\date{17 October 2026}
\setlength{\emergencystretch}{5em}
\begin{document}
\maketitle
\begin{abstract}
We derive an exact local, gauge-covariant representation
of the actual leading determinant using zero covariant
block means in the scalar gauge. Both its scalar and
vector factors have flat spectral gaps of order
\(K^{-2}\), with uniform heat and inverse-power bounds.
The finite amplitude identity continues jointly through
the flat stabilizers and passes to local infinite
systems at each fixed \(K\). Its invariant diagonal
densities give the actual Maxwell moment directly.
This closes the specific infrared passage left open
by the earlier massless physical quotient.

The exact finite physical quotient of Notebook 34 requires
joint control of its scalar and physical determinants.
We give an explicit reason: a globally based massless
scalar determinant on a periodic constant Cartan connection
has a nonzero quadratic response per volume, whereas the
ordered thermodynamic and zero-mass limit gives zero.
The fine/coarse scalar ratio has an extra term proportional
to \(K^2-1\). The actual local gauge Ward identity and
the finite quotient formula prove that the physical
determinant cancels this term exactly in the thermodynamic
limit. No flux seam or nonperiodic magnetic family is used.

We also extend the flat physical spectral estimates to
uniform heat-trace and inverse-power bounds. The
rescaled inverse square has at most logarithmic trace
growth; powers above two have uniformly bounded trace
per final block. A local tent lift proves full
Laplacian domination for the new vector gauge, so
both new factors have the corresponding estimates.
Their endpoint-uniform curvature comparison remains
the next step toward a positive actual drift.
\end{abstract}

\section*{Scope and conventions}

All actual operators are those of the homogeneous local
construction in Notebooks 32--34. The primitive blocking
ratio is fixed, \(K=\ell^k\), and \(M\) denotes the
preceding homogeneous depth. Traces count real adjoint
coordinates. For a normalized Cartan element \(H\),
\(\tau(H^2)=1\), set
\[
 C_H=\sum_{q\in\operatorname{spec}(\operatorname{ad}_H)}q^2,
 \qquad \alpha_g=\frac{C_H}{48\pi^2}.
\]
The complexified charge decomposition counts each real
adjoint degree of freedom once in this trace convention.
The notation \([t^2]\) denotes a Taylor coefficient,
not a second derivative.

The constant flat connection used below tests the
zero-momentum Ward identity. It is different from the
constant-curvature probe that extracts the Maxwell
coefficient. The infrared order of limits is stated
explicitly for each result.



% No preamble. Prefix jsi-. Ordered scalar limit and exact holonomy obstruction.
\section{The scalar infrared limit and a forced holonomy cancellation}
\label{jsi-section}

There are two different assertions in this section. A common positive
physical mass gives a well-defined thermodynamic scalar reference,
whose subsequent massless limit is \(\alpha_g\log K\). The original
separately massless globally based determinants need not have that
order-independent analytic limit. An exactly periodic flat connection
gives a concrete obstruction and, through the actual finite determinant
identity, forces a compensating response in the physical factor.

\subsection{The mass convention and an exact base correction}
\label{jsi-mass-section}

Fix \(K=\ell^k\) and final torus side \(N\) in four dimensions.
Current fine sites have physical spacing \(K^{-1}\), with side \(KN\);
coarse sites have spacing one and side \(N\).
All color coordinates are trace orthonormal.
Let \(G_{\rm all}(U)\) be the covariant scalar Laplacian on all sites,
and let \(G_b(U)\) be its principal minor after deleting the color
coordinates at one common physical base \(o\). Deleting a site does
not change the remaining diagonal entries.

A common physical mass \(\delta>0\) means
\begin{equation}
 H_{f,\delta}=K^2G_{{\rm all},f}+\delta^2I,\qquad
 H_{c,\delta}=G_{{\rm all},c}+\delta^2I.
 \label{jsi-mass}
\end{equation}
Their principal minors are \(H_{f,\delta}^b,H_{c,\delta}^b\).
With \(\mathcal N_0F=F(W)-F(1)\), define
\begin{equation}
 S_{\delta,N}^b=\frac12\mathcal N_0
       \bigl(\log\det H_{f,\delta}^b-\log\det H_{c,\delta}^b\bigr).
 \label{jsi-scalar-ratio}
\end{equation}
The fine counting mass is \(\delta/K\), not \(\delta\).
At zero mass multiplication of the fine operator by \(K^2\)
changes only a background-independent normalization.
Equation \eqref{jsi-scalar-ratio} is a specified scalar comparison
family. It is not automatically an exact mass deformation of the
entire physical-quotient determinant identity.

\begin{lemma}[Positive-mass base defect]
\label{jsi-base-lemma}
Put \(H=\epsilon^{-2}G_{\rm all}+\delta^2I\), \(R=H^{-1}\),
where \(\epsilon=K^{-1}\) or \(1\). Then
\begin{equation}
       \det H^b=\det H\,\det R_{oo}.
 \label{jsi-jacobi}
\end{equation}
At fixed \(\delta>0,K\), the normalized logarithm of the fixed color
block \(R_{oo}\) has a complex product neighborhood and first-two
bounded-direction derivative bounds independent of \(N\).
Its first-two source kernels are exponentially summable in their
distances from \(o\), including all fixed polynomial moments.
No uniformity as \(\delta\downarrow0\) is asserted.
\end{lemma}

\begin{proof}
The determinant identity is the block Jacobi identity. At real
backgrounds,
\[
 \delta^2I\le H\le(16\epsilon^{-2}+\delta^2)I,
 \qquad
 (16\epsilon^{-2}+\delta^2)^{-1}I
                      \le R_{oo}\le\delta^{-2}I.
\]
The nearest-neighbor coefficients have volume-independent operator
bounds on a small complex product ball. The resolvent identity,
and then the lower bound for \(R_{oo}\), give a smaller complex
ball on which both inverses remain bounded uniformly in volume.
The logarithm branch continued from zero is bounded there after
normalization. Its first two derivatives also follow directly from
\begin{align}
 DR[v]&=-RH'[v]R,\notag\\
 D^2R[v,w]&=RH'[v]RH'[w]R+RH'[w]RH'[v]R-RH''[v,w]R.
 \label{jsi-resolvent-derivatives}
\end{align}
Differentiating \(\log\det R_{oo}\) uses only its fixed color trace.

For exponential bounds, conjugate by a site weight whose ratios
across nearest-neighbor edges differ from one by \(O(\kappa)\).
Choose \(\kappa>0\), depending on \(\delta,\epsilon\), so the
conjugation perturbation has norm at most \(\delta^2/4\).
The weighted inverse is then bounded by \(2/\delta^2\).
Spending a smaller exponent gives absolute row and column Schur
bounds. The derivatives \(H',H''\) are local. Every source chain
in \eqref{jsi-resolvent-derivatives}, and in the logarithmic
derivatives, begins and ends at \(o\). Its exponential weights
therefore control all source distances from \(o\), with a further
exponent reserve for any fixed polynomial moment.
\end{proof}

Writing \(S_{\delta,N}^{\rm all}\) for the unbased half-ratio gives
the exact comparison
\begin{equation}
 S_{\delta,N}^b=S_{\delta,N}^{\rm all}
 +\frac12\mathcal N_0\bigl\{
       \log\det(R_f)_{oo}-\log\det(R_c)_{oo}\bigr\}.
 \label{jsi-defect}
\end{equation}
The lemma has two useful consequences. A base receding from fixed
localized sources changes their first-two coefficients
exponentially little. Alternatively, average the based functional
over the \(N^4\) simultaneous coarse translations of the fine and
coarse base. Reanchoring its Hessian at one output cell bounds the
anchored second-moment norm of its difference from the unbased
functional by
\begin{equation}
                         C_{\delta,K}N^{-4}.
 \label{jsi-average-defect}
\end{equation}
Indeed summing the translated rooted defect kernel converts that
row into its full source sum, divided by \(N^4\); its displacement
moments are bounded by the lemma. This averaging is an estimate
for the limit, not a claim that different finite base choices
give identical functionals.

After composition with the actual fixed-\(K\) stationary branch,
the same conclusions use its previously proved localized source
kernels. Under the Notebook 33 inputs the constants are uniform
in old depth \(M\), still depending on \(K,\delta\).

\subsection{The ordered localized Maxwell limit}
\label{jsi-ordered-section}

At positive mass, the unbased diagonal logarithm has exponentially
localized Taylor kernels, by the same resolvent argument applied
to its logarithm integral. Its color trace is gauge invariant.
Its Hessian moments therefore define its affine curvature
coefficient absolutely. This gives a way to take the limit without
postulating an affine magnetic field on a torus: first pass to
infinite kernels using compactly supported directions, then
evaluate polynomial directions by their exponential moments.
Symmetric short lifts of periodic kernels give the same result.
At fixed \(K,\delta\), unbased winding contributions vanish
exponentially in \(N\), and \eqref{jsi-average-defect} removes the
averaged base contribution.

For an actual stationary insertion, use fixed-\(K\) structural
compactness followed by the polynomial-jet identification of
Notebook 33. Every such limit has the same scalar expression
below. This statement uses the already established source kernels;
it does not infer locality from a finite determinant identity.

Normalize the Maxwell functional by
\(\beta(F)=2[t^2]\) of its localized constant-curvature density.
For \(\mu\ne\nu\) and \(\tau(H^2)=1\), the current and coarse
Cartan probes are
\[
 U_\mu(t;x)=e^{itx_\nu H/K^2},\qquad
 W_\mu(t;y)=e^{ity_\nu H}.
\]
The unpinned scalar heat coefficient already established is
\begin{equation}
 \begin{split}
 h(s)&=\frac{C_H}{3}s^2e^{-8s}I_0(2s)^2I_1(2s)^2,\\
 h(s)&=O(s^4)\quad(s\downarrow0),\\
 h(s)&=\alpha_g+O(s^{-1})\quad(s\to\infty),
 \qquad \alpha_g=\frac{C_H}{48\pi^2}.
 \end{split}
 \label{jsi-heat}
\end{equation}
The \(K^4\) fine sites per final block cancel the curvature
factor \(K^{-4}\). In common physical heat time \(u\), the fine
and coarse coefficients are consequently \(h(K^2u)\) and \(h(u)\).

\begin{proposition}[Ordered scalar reference limit]
\label{jsi-ordered-theorem}
Take the thermodynamic localized-kernel limit first, at fixed
\(\delta>0\) and fixed \(K\). The resulting scalar coefficient is
\begin{equation}
 \beta_\delta(K)=\frac12\int_0^\infty
           e^{-\delta^2u}\{h(K^2u)-h(u)\}\frac{du}{u}.
 \label{jsi-beta-delta}
\end{equation}
It satisfies
\begin{equation}
             \lim_{\delta\downarrow0}\beta_\delta(K)
                         =\alpha_g\log K.
 \label{jsi-beta-limit}
\end{equation}
For \(K\ge1\) and \(0<\delta\le1/2\), the explicit reference
also obeys
\begin{equation}
 |\beta_\delta(K)-\alpha_g\log K|
       \le C\delta^2\{1+\log(1/\delta)\},
 \label{jsi-mass-rate}
\end{equation}
where \(C\) is independent of \(K\).
\end{proposition}

\begin{proof}
The localized logarithm heat formula, the preceding normalization,
and \eqref{jsi-heat} give \eqref{jsi-beta-delta}.
Without the mass factor, the paired integrand is absolutely
integrable: use its small-\(s\) bound at zero and cancellation
of \(\alpha_g\) at infinity. On \(\varepsilon<u<T\), its
unhalved integral equals
\[
 \int_T^{K^2T}h(v)\frac{dv}{v}
       -\int_\varepsilon^{K^2\varepsilon}h(v)\frac{dv}{v}.
\]
The limits are \(2\alpha_g\log K\) and zero, respectively.
Dominated convergence proves \eqref{jsi-beta-limit}.

For the rate, \(h\) is bounded on the whole positive axis.
On \(0<u<1\), use \(1-e^{-\delta^2u}\le\delta^2u\);
the resulting error is \(O(\delta^2)\), uniformly in \(K\).
For \(u\ge1\), \eqref{jsi-heat} gives
\(|h(K^2u)-h(u)|\le2C/u\) for every \(K\ge1\).
Split the remaining integral at \(u=\delta^{-2}\).
The first part is \(O(\delta^2\log(1/\delta))\), the
second \(O(\delta^2)\).
\end{proof}

The uniformity in \eqref{jsi-mass-rate} belongs to the explicit
reference after the thermodynamic limit. It does not interchange
mass removal with that limit for the original based determinants.

\subsection{An exact massless anomaly on a periodic flat connection}
\label{jsi-anomaly-section}

There is a direct obstruction to such an interchange, already
at zero momentum. Consider one charged scalar \(q\) on a periodic
\(n^4\) lattice, with constant phase \(e^{itqA}\) on direction
\(\mu\). This is a legitimate periodic flat connection for every
\(t\); no continuously varying magnetic flux is involved.
Its unbased eigenvalues are
\begin{equation}
 \lambda_p(t)=2\sum_\nu
       \{1-\cos(p_\nu+tqA\,\delta_{\mu\nu})\},
       \qquad p\in(2\pi/n)\mathbb Z_n^4.
 \label{jsi-spectrum}
\end{equation}
Put \(V=n^4\), and delete one scalar site to obtain \(G_n^b(t)\).
Translation invariance of the resolvent diagonal and
\eqref{jsi-jacobi} imply
\begin{equation}
 \det G_n^b(t)=\frac1V\prod_{p\ne0}\lambda_p(t)
       \left\{1+\lambda_0(t)
                    \sum_{p\ne0}\lambda_p(t)^{-1}\right\}.
 \label{jsi-cofactor}
\end{equation}
Initially this is an identity away from zero eigenvalues;
both sides continue analytically through \(t=0\).

\begin{proposition}[Noncommuting mass and thermodynamic derivatives]
\label{jsi-anomaly-theorem}
Let
\begin{equation}
 g_4=\int_{[-\pi,\pi]^4}\frac{dp}{(2\pi)^4}
             \frac1{2\sum_\nu(1-\cos p_\nu)}>0.
 \label{jsi-green}
\end{equation}
Then
\begin{equation}
 \lim_{n\to\infty}n^{-4}[t^2]\log\det G_n^b(t)
                           =q^2A^2g_4.
 \label{jsi-anomaly}
\end{equation}
At every fixed positive scalar mass, the corresponding
thermodynamic coefficient is zero, and remains zero when
the mass is subsequently removed.
\end{proposition}

\begin{proof}
Since \(\lambda_0(t)=q^2A^2t^2+O(t^4)\),
\eqref{jsi-cofactor} gives, writing \(\lambda_p=\lambda_p(0)\),
\begin{equation}
 [t^2]\log\det G_n^b(t)
   =q^2A^2\sum_{p\ne0}
    \left\{\frac{\cos p_\mu}{\lambda_p}
       -\frac{2\sin^2p_\mu}{\lambda_p^2}
       +\frac1{\lambda_p}\right\}.
 \label{jsi-cofactor-jet}
\end{equation}
The first two summands equal
\(\partial_{p_\mu}(\sin p_\mu/\lambda_p)\).
Their singularity is \(O(|p|^{-2})\), integrable in dimension
four. Removing a radius-\(\varepsilon\) ball before integration
by parts leaves boundary flux \(O(\varepsilon^2)\), so their
Brillouin-zone integral is zero. Their grid sums converge:
\begin{equation}
 n^{-4}\sum_{0<|p|\le\varepsilon}|p|^{-2}
                       \le C\varepsilon^2+O(n^{-2}).
 \label{jsi-grid-bound}
\end{equation}
Outside that ball ordinary Riemann convergence applies.
The same argument for \(\lambda_p^{-1}\) gives \(g_4\),
proving \eqref{jsi-anomaly}.

At fixed positive mass \(m\), the unbased coefficient per volume
instead tends to the integral of
\(\partial_{p_\mu}\{\sin p_\mu/(\lambda_p+m^2)\}\),
which is zero by periodicity. Lemma~\ref{jsi-base-lemma}
makes the based correction \(O(1)\) in these derivatives,
so division by \(n^4\) removes it. Subsequent \(m\downarrow0\)
still gives zero.
\end{proof}

The deleted block has fixed rank, but its massless derivative
contribution is extensive. Equation \eqref{jsi-cofactor} identifies
the soft global mode responsible for the failure. Summing
adjoint charges replaces \(q^2\) by \(C_H\); for a vector of
commuting constant angles the answer is \(C_H|A|^2g_4\),
with mixed terms zero by coordinate reflections.

\subsection{The actual branch and its forced physical compensation}
\label{jsi-actual-section}

Now use the actual final constant links and current representative
\begin{equation}
 W_\mu(t)=e^{itA_\mu H},\qquad
 U_\mu(t)=e^{itA_\mu H/K}.
 \label{jsi-flat-branch}
\end{equation}
Each ordinary blocking stage preserves this commuting connection:
all paths with the same displacement have the same transporter,
so the selected mean is evaluated on identical entries.
The same flat lift exists through every old refinement \(M\).
All plaquettes are identity. Since the real Wilson action is
nonnegative, the associated classical minimum values are zero.
After the actual local root gauge, uniqueness of the small
minimum germ identifies this with the actual stationary branch.
This is a fixed finite-germ argument.

Every sitewise adjoint gauge preserves a scalar zero at the
global base. Thus the globally based scalar determinants are
unchanged by that root gauge and may be computed in
\eqref{jsi-flat-branch}.
Fine current volume is \((KN)^4\) and fine angle is \(tA/K\);
coarse volume is \(N^4\) and angle is \(tA\).
Proposition~\ref{jsi-anomaly-theorem} therefore gives
\begin{equation}
 \lim_{N\to\infty}N^{-4}[t^2]
         (\log\det G_f-\log\det G_c)
             =C_H|A|^2(K^2-1)g_4.
 \label{jsi-orbit-anomaly}
\end{equation}
The generators on this branch are the same for every old \(M\).
This is a flat-holonomy, zero-momentum response, not a Maxwell
coefficient.

\begin{proposition}[Forced compensation in the physical factor]
\label{jsi-compensation-theorem}
Fix \(K\) before taking the thermodynamic limit.
For the actual leading correction and physical quotient,
\begin{align}
 \lim_{N\to\infty}N^{-4}[t^2]D_N(W(t))&=0,
       \label{jsi-total-zero}\\
 \lim_{N\to\infty}N^{-4}[t^2]\log\det Q_\perp(W(t))
                  &=C_H|A|^2(K^2-1)g_4.
       \label{jsi-compensation}
\end{align}
The same conclusion allows varying old depth under the
previously established uniform-old-depth, fixed-\(K\) local
kernel bounds.
\end{proposition}

\begin{proof}
The actual correction \(D_N\) has translation-covariant Hessian
kernels with common exponential row and source bounds at fixed
\(K\). Every structural limit has invariant local densities and
zero linear color jet. Its quadratic Ward identity consequently
annihilates lattice gradients. Exponential moments extend this
identity to linear scalar gauge parameters by cutoff approximation.
Their gradients are constant link fields, so the infinite
quadratic symbol vanishes at momentum zero.

The finite normalized coefficient on \eqref{jsi-flat-branch} is
one-half the contracted Hessian row sum. Uniform exponential
tails make its distant part uniformly small. Every fixed
central sum converges along a structural subsequence, and its
full limiting row sum is zero by the preceding Ward argument.
All subsequential scalar limits are therefore zero, proving
\eqref{jsi-total-zero}. The uniform-old-\(M\) versions of these
same bounds permit \(M\) to vary. No uniform-in-\(K\) locality
is used.

At finite \(N\) the exact physical-quotient identity is
\begin{equation}
 2D_N=\mathcal N_0
       (\log\det G_f-\log\det G_c-\log\det Q_\perp).
 \label{jsi-actual-identity}
\end{equation}
Subtract \eqref{jsi-total-zero} from
\eqref{jsi-orbit-anomaly} using the same branch, base and
regulator. This yields \eqref{jsi-compensation}, including
its sign and normalization.
\end{proof}

No claim that \(D_N(W(t))\) is identically zero at finite \(N\)
is needed: a periodic constant connection may have nontrivial
holonomy. Nevertheless, its thermodynamic constant-field
response must vanish by locality and the Ward identity.
Thus the term \(-\frac12\log\det Q_\perp\) cancels the scalar
half-ratio exactly at zero momentum.

\subsection{What the ordered limit does not identify}
\label{jsi-scope}

The positive result is
\[
 \lim_{\delta\downarrow0}\lim_{N\to\infty}
       \{\hbox{localized regulated scalar Maxwell coefficient}\}
                         =\alpha_g\log K.
\]
The original separately massless globally based family does not
have an order-independent analytic local limit, as the periodic
test \eqref{jsi-orbit-anomaly} proves.
In particular a \(K\)-uniform \(O(1)\) constant-potential response
for the separate physical determinant is false, despite its
flat lower bound \(cK^{-2}\).

The anomaly alone neither computes nor disproves equality of
Maxwell coefficients after a suitable coupled cancellation.
A valid joint vector/scalar regulator must preserve the
rank-fixed but derivative-large base factors before its
infrared limit. The exact finite quotient identity does not
provide that regulated identity by itself. No continuum
measure, nontriviality, or mass-gap conclusion follows here.


% No preamble. Prefix pvs-. Actual flat physical spectral estimates.
\section{Flat physical heat densities and the remaining curvature interface}
\label{pvs-section}

Fix dimension four, a color algebra of dimension \(n_{\mathfrak g}\),
a primitive odd ratio \(\ell\), a composite \(K=\ell^k\), and any
preceding homogeneous depth \(M\).
Coordinates are current-link angles with counting inner products.
On a cubic current torus of side \(KN\), there are \(N^4\) final
\(K\)-blocks.
All statements about varying backgrounds below are explicitly
separated from the flat spectral theorem.

\subsection{Actual flat inputs}
\label{pvs-input-section}

Let \(P_f,P_c\) be the fine and coarse physical projections at
the flat field, and write
\begin{equation}
 \begin{split}
 L&=q_M,\qquad A=P_cC_K|_{\operatorname{ran}P_f},\\
 R&=L|_{\operatorname{ran}P_f}-A^Tq_{MK}A,\qquad
 Q=R+aA^TA .
 \end{split}
 \label{pvs-operators}
\end{equation}
These are the actual residual physical operators of the finite
quotient identity. The following previously proved facts are the
only model-specific inputs:
\begin{equation}
 \begin{gathered}
 L\ge d^Td,\qquad 0\le q_{MK},\qquad \|q_{MK}\|\le R_q,\\
 R\ge0,\qquad Q\ge c_aK^{-2}I,\qquad c_a>0 .
 \end{gathered}
 \label{pvs-inputs}
\end{equation}
The constants are uniform in \(M,K,N\).
For example, the physical alias proof gives
\[
 c_a=\min\{B_0^{-1},a/(A_0+B_0R_q)\},\qquad
 A_0=2c_\omega^{-2},\quad
 B_0=(1+2c_\omega^{-2})/4,\quad c_\omega=(2/\pi)^5.
\]
Flat operators are invariant under final-block translations.

The globally based scalar gradient has the same range as the
ordinary periodic gradient: subtracting the value of a scalar at
the base does not change its gradient.
Thus \(\operatorname{ran}P_f\) is the usual transverse Bloch
space, including all harmonic one-forms at momentum zero.
No estimate on an inverse globally based scalar operator is used.

Choose a fixed \(a\ge R_q\). Then
\begin{equation}
 Q=L|_{\rm phys}+A^T(aI-q_{MK})A\ge\Delta_{\rm phys},
 \qquad
 Q\ge\tfrac12\Delta_{\rm phys}+\frac{c_a}{2K^2}I .
 \label{pvs-combined-order}
\end{equation}
Here \(\Delta_{\rm phys}\) is the Maxwell form on the transverse
space, where it is the scalar lattice Laplacian.

\begin{theorem}[Uniform flat physical heat and inverse moments]
\label{pvs-main}
Put \(\widehat Q=K^2Q\), and define the heat trace per final block by
\[
 \widehat H_{M,K,N}(u)
   =N^{-4}\operatorname{Tr}_{\rm phys}e^{-u\widehat Q}.
\]
There are \(C,c>0\), independent of \(M,K,N\), such that
\begin{equation}
 \widehat H_{M,K,N}(u)
 \le C\left\{\frac{K^4}{(1+K^2u)^2}+N^{-4}\right\}e^{-cu},
 \qquad u>0.
 \label{pvs-heat}
\end{equation}
In the infinite flat structural system the corresponding block
trace satisfies the same bound with \(N^{-4}\) omitted.
For every \(p>0\),
\begin{equation}
 \operatorname{Tr}_{\rm block}\widehat Q^{-p}
 \le C_p
 \begin{cases}
 K^{4-2p}+N^{-4},&0<p<2,\\
 1+\log K+N^{-4},&p=2,\\
 1+N^{-4},&p>2,
 \end{cases}
 \label{pvs-moments}
\end{equation}
again omitting \(N^{-4}\) in infinite volume.
For each fixed real \(r\) and \(U\ge1\),
\begin{equation}
 \int_U^\infty u^{r-1}\widehat H_{M,K,N}(u)\,du
 \le C_r(1+N^{-4})e^{-cU/2}.
 \label{pvs-tail}
\end{equation}
These are unperturbed spectral statements. They do not assert a
bound on a background derivative of a heat trace.
\end{theorem}

\begin{proof}
For self-adjoint finite matrices \(X\ge Y\), the min--max theorem
gives ordered eigenvalues \(\lambda_j(X)\ge\lambda_j(Y)\);
therefore \(\operatorname{Tr}e^{-sX}\le
\operatorname{Tr}e^{-sY}\).
No operator monotonicity of the exponential is asserted.

At momentum \(p\in(2\pi/n)\mathbb Z_n^4\), \(n=KN\), put
\(P(p)=4\sum_\mu\sin^2(p_\mu/2)\).
The physical Maxwell multiplicity is \(3n_{\mathfrak g}\)
away from zero and \(4n_{\mathfrak g}\) at zero. Hence
\begin{equation}
 \operatorname{Tr}_{\rm phys}e^{-sQ}
 \le 4n_{\mathfrak g}e^{-c_as/(2K^2)}
                       \sum_p e^{-sP(p)/2}.
 \label{pvs-minmax}
\end{equation}
For the one-dimensional scalar heat average
\[
 p_n(t)=\frac1n\sum_{m\ {\rm mod}\ n}
                      e^{-4t\sin^2(\pi m/n)},
\]
use representatives \(|m|\le n/2\) and
\(\sin(\pi|m|/n)\ge2|m|/n\).
Gaussian summation for \(t\ge1\), and \(p_n(t)\le1\) for
\(t\le1\), give
\begin{equation}
 p_n(t)\le n^{-1}+C(1+t)^{-1/2},\qquad
 p_n(t)^4\le C\{n^{-4}+(1+t)^{-2}\}.
 \label{pvs-scalar-heat}
\end{equation}
Since \(\sum_p e^{-tP(p)}=n^4p_n(t)^4\),
division of \eqref{pvs-minmax} by \(N^4\) yields
\[
 N^{-4}\operatorname{Tr}e^{-sQ}
 \le C\{K^4(1+s)^{-2}+N^{-4}\}e^{-cs/K^2}.
\]
Set \(s=K^2u\) to obtain \eqref{pvs-heat}.

The normalized total trace equals the diagonal trace in any one
final block, by final-block translation invariance.
This equality does not turn the min--max argument into a
pointwise operator comparison.
In infinite volume, apply the same min--max argument in each
finite-dimensional coarse Bloch fiber and integrate.
Folding the Maxwell aliases gives \(K^4\) times the fine scalar
heat density, proving the infinite statement directly.

Positive spectral calculus gives
\[
 \operatorname{Tr}_{\rm block}\widehat Q^{-p}
  =\frac1{\Gamma(p)}\int_0^\infty
      u^{p-1}\operatorname{Tr}_{\rm block}e^{-u\widehat Q}\,du .
\]
Split at \(K^{-2}\) and \(1\) to obtain
\eqref{pvs-moments}.
For \(p>2\), one can instead use
\(K^4/(1+K^2u)^2\le u^{-2}\) directly.
The same inequality and exponential decay prove
\eqref{pvs-tail}.
\end{proof}

The finite \(N^{-4}\) term is necessary: it retains the finite
number of physical harmonic modes.
The local inverse moment means
\(\operatorname{Tr}(P_{\rm block}\widehat Q^{-p}P_{\rm block})\).
It is not a claim of global trace class in infinite volume, nor
an identification with a power of the compressed inverse.
The logarithm in the \(p=2\) bound is an upper bound; no matching
actual asymptotic is claimed.

\subsection{An optional finite on-shell replacement}
\label{pvs-onshell-section}

On a finite real on-shell family with
\(LD_f=0\) and \(qD_c=0\), set
\[
 Q_{\rm on}=L|_{\rm phys}+aA^TA=Q+A^Tq_{\rm phys}A.
\]
The minimum right inverse gives \(AQ^{-1}A^T=a^{-1}I\).
Consequently the determinant lemma proves
\begin{equation}
 \frac{\det Q_{\rm on}}{\det Q}
   =\det(I+a^{-1}q_{\rm phys}).
 \label{pvs-onshell-determinant}
\end{equation}
Because \(q\) kills the coarse gauge space, the normalized
logarithm on the right is the normalized full coarse
\(\log\det(q+aI)\); only a field-independent power of \(a\)
has been removed.
The already proved coarse-factor Maxwell estimate can therefore
control this replacement wherever its Maxwell moment is
legitimately defined.

This does not apply to an arbitrary off-shell family by omission
of the multiplier curvature.
The actual general form is
\(L=s''-\lambda\cdot f''\).
An infinite affine use of separate globally based physical
determinants requires its own infrared justification.

\subsection{A precise sufficient derivative estimate}
\label{pvs-vertex-section}

The following fixed-space lemma isolates an additional estimate
which would suffice for an infrared second-variation bound.
Its hypotheses are not established here for the actual varying
physical quotient.

\begin{lemma}[Relative-form vertices]
\label{pvs-vertex-lemma}
Let
\(X(t)=X+tV_1+t^2V_2+O(t^3)\)
be a real self-adjoint family on a fixed finite-dimensional
Hilbert space, with \(X>0\). Suppose
\begin{equation}
 \|X^{-1/2}V_jX^{-1/2}\|\le L_j,\qquad j=1,2.
 \label{pvs-vertex-hyp}
\end{equation}
Then
\begin{equation}
 \left|[t^2]\operatorname{Tr}e^{-uX(t)}\right|
 \le8(L_2+L_1^2)\operatorname{Tr}e^{-uX/4},
 \qquad u>0.
 \label{pvs-vertex-bound}
\end{equation}
\end{lemma}

\begin{proof}
The single insertion has absolute value at most
\[
 uL_2\operatorname{Tr}(Xe^{-uX})
 \le(2/e)L_2\operatorname{Tr}e^{-uX/2}.
\]
In an eigenbasis of \(X\), cyclic Duhamel gives the double
insertion exactly as
\begin{equation}
 \frac u2\sum_{i,j}|(V_1)_{ij}|^2
       \frac{e^{-u\lambda_i}-e^{-u\lambda_j}}
                         {\lambda_j-\lambda_i},
 \label{pvs-divided-difference}
\end{equation}
with the continuous value \(ue^{-u\lambda_i}\) at coincident
eigenvalues.
For \(a,b>0\),
\begin{equation}
 ab\,\frac{e^{-ua}-e^{-ub}}{b-a}
 \le2\{ae^{-ua/2}+be^{-ub/2}\}.
 \label{pvs-elementary}
\end{equation}
Indeed, assume \(b\ge a\).
For \(b\ge2a\), use \(ab/(b-a)\le2a\).
For \(a\le b\le2a\), use the derivative bound
\(ue^{-ua}\), \(ab\le2a^2\), and
\((ua)e^{-ua/2}\le2/e\).

Write \(V_1=X^{1/2}BX^{1/2}\), with \(\|B\|\le L_1\).
The row and column square sums of \(B\) are at most \(L_1^2\).
Applying \eqref{pvs-elementary} to
\eqref{pvs-divided-difference} bounds it by
\[
 2uL_1^2\operatorname{Tr}(Xe^{-uX/2})
 \le(8/e)L_1^2\operatorname{Tr}e^{-uX/4}.
\]
Together with the single insertion this proves the stated
conservative constant.
\end{proof}

If \(X(t)=K^2Q_\perp(t)\) could be identified with a fixed
physical space and satisfy \eqref{pvs-vertex-hyp} uniformly,
Theorem~\ref{pvs-main} would give
\begin{equation}
 \int_1^\infty
 \left|N^{-4}[t^2]\operatorname{Tr}e^{-uX(t)}\right|
 \frac{du}{u}
 \le C(L_2+L_1^2)(1+N^{-4}).
 \label{pvs-conditional-tail}
\end{equation}
An infinite-volume version additionally requires a justified
trace-and-derivative limit.
Neither the uniform identification of varying physical spaces
nor the actual bounds \eqref{pvs-vertex-hyp} follows from the
flat gap.
Derivatives of \(P_f,P_c\), the current, and the nonlinear
constraint cannot be absorbed into an unproved constant.

\subsection{The required scalar--physical cancellation}
\label{pvs-scope-section}

The separate physical determinant has a growing response even
along constant commuting final link angles.
The actual globally based finite identity and the local Ward
limit give, at fixed \(K\),
\begin{equation}
 \lim_{N\to\infty}N^{-4}[t^2]\log\det Q_\perp(tA)
   =C_H|A|^2(K^2-1)g_4,\qquad
 g_4=\int_{[-\pi,\pi]^4}
       \frac{dp}{(2\pi)^4\,4\sum_\mu\sin^2(p_\mu/2)}>0 .
 \label{pvs-holonomy}
\end{equation}
Here the fine angles are \(tA/K\).
This is the separately established flat-holonomy calculation:
the same term in the scalar orbit-density ratio cancels
\eqref{pvs-holonomy} in the actual total correction.
It is a zero-momentum test, not a magnetic Maxwell coefficient.

Accordingly, the spectral theorem does not permit an independent
replacement of a massless physical or scalar factor by its
massive infinite reference.
A valid curvature comparison must preserve the coupled
cancellation, or supply a justified joint infrared regulator.
No actual vector logarithmic coefficient, its sign, or
Yang--Mills existence conclusion is proved by this section.



% No preamble. Prefix ccs-. Actual covariant moving-reference tangents.
\section{A covariant moving-reference scalar average}
\label{ccs-section}

We construct tangent operators for the actual ordinary composite
average. The nonlinear constraint and its cutoff are unchanged.
The construction uses the internal group-valued transport selected
in that ordinary map, rather than an unverified identification with
an older axial operator carrying the name \(Q'\).

\subsection{Actual data and the primitive relative correction}
\label{ccs-primitive-section}

Fix odd \(\ell\ge3\), \(K=\ell^k\), and compatible nested centered
blocks. Coordinates are current-link angles with counting scalar
and bond inner products. Use
\[
 U^g(e)=g(e_-)U(e)g(e_+)^{-1},\qquad
 (D_U\phi)(e)=\phi(e_-)-\operatorname{Ad}_{U(e)}\phi(e_+).
\]
The generator acts on all scalar sites. Scalar parameters have
angle units at every level; no spacing factor is inserted into
this formula.

Let \(P_\ell(U;y,x)\) be exactly the internal group-valued transport
entering the selected ordinary map \(M_\ell\). The actual definition
in Ba{\l}aban, CMP109 (1987), pp.~252--254, (0.3),(0.5)--(0.12),
supplies local analyticity, internal-block support, and
\begin{equation}
 P_\ell(U^g;y,x)=g(y)P_\ell(U;y,x)g(x)^{-1}.
 \label{ccs-endpoint}
\end{equation}
At identity its derivative is the same arithmetic contour mean
that occurs in the ordinary arithmetic differential.
These facts also hold when \(P_\ell\) is the selected group mean
of coordinate-order transporters; it need not be one axial path.

The pinned deposition specifies this ordinary map at Raw I:
831--845. Its relative contour and ordinary increments occur at
46965--46978, and the moving-reference endpoint correction at
47076--47086. Raw I:23913--23920 describes an older axial/nested
scalar average. No equality with every such convention is needed:
the following operators are defined directly using the actual
\(P_\ell\).

Freeze \(U\), put \(W=M_\ell(U)\), and for a small bond increment
\(x\) define
\begin{equation}
 \begin{split}
 r_U(x;y,z)
   &=\frac1i\log\!\left[
       P_\ell(e^{ix}U;y,z)P_\ell(U;y,z)^{-1}\right],\\
 \alpha_U(x;y)&=\ell^{-d}\sum_{z\in B_\ell(y)}r_U(x;y,z),\\
 a_U(x;y)&=e^{i\alpha_U(x;y)}.
 \end{split}
 \label{ccs-relative-mean}
\end{equation}
All relative logarithms vanish at \(x=0\), and \(a_U(0)=1\).
Let
\begin{equation}
 Z_Ux=D\alpha_U(0)[x],\qquad
 C_Ux=\frac1i\left.\frac d{ds}\right|_0
                  M_\ell(e^{isx}U)W^{-1}.
 \label{ccs-tangents}
\end{equation}
The local corrected output with this frozen reference is
\[
 \overline F_U(x;c)=a_U(x;c_-)^{-1}
                         M_\ell(e^{ix}U)(c)a_U(x;c_+).
\]
It equals \(W\) at zero. Its left-relative differential is exactly
\begin{equation}
                     \overline C_U=C_U-D_WZ_U.
 \label{ccs-primitive-correction}
\end{equation}
This is a moving-reference tangent identity. It is not an
identification with the derivative of the fixed-identity corrected
map at \(U\).

\begin{lemma}[Primitive scalar intertwining]
\label{ccs-primitive-lemma}
Define
\begin{equation}
 (B_U\phi)(y)=\ell^{-d}\sum_{z\in B_\ell(y)}
                  \operatorname{Ad}_{P_\ell(U;y,z)}\phi(z).
 \label{ccs-primitive-B}
\end{equation}
With \(r_\ell\) restriction to the new centers,
\begin{equation}
 Z_UD_U=r_\ell-B_U,\qquad
                  \overline C_UD_U=D_WB_U.
 \label{ccs-primitive-intertwining}
\end{equation}
\end{lemma}

\begin{proof}
For the fine gauge \(e^{is\phi}\), the relative contour derivative
from \eqref{ccs-endpoint} is
\[
 \phi(y)-\operatorname{Ad}_{P_\ell(U;y,z)}\phi(z).
\]
Its fine left tangent is \(D_U\phi\). Averaging this expression
proves the first identity. There is no logarithm Jacobian at this
step: its argument is identity at the reference point.
Ordinary covariance gives \(C_UD_U=D_Wr_\ell\). Subtract
\(D_WZ_UD_U\) using \eqref{ccs-primitive-correction}.
\end{proof}

\subsection{Joint covariance and composition}
\label{ccs-composition-section}

For a fixed site gauge \(g\), let \(O_E(g),O_S(g),O_c(g)\)
denote adjoint frames on fine bond source sites, fine scalar sites,
and new centers. Since
\[
 e^{iO_E(g)x}U^g=(e^{ix}U)^g,
\]
the relative contour product at \(U^g\) is conjugated by \(g(y)\).
The similarity-equivariant logarithm therefore gives
\begin{equation}
 \begin{split}
 Z_{U^g}O_E(g)&=O_c(g)Z_U,\\
 \overline C_{U^g}O_E(g)&=O_{E_c}(g)\overline C_U,\\
 B_{U^g}O_S(g)&=O_c(g)B_U.
 \end{split}
 \label{ccs-primitive-covariance}
\end{equation}
Both the reference and the varied configuration are transformed.
No multiplicative cocycle in the reference is assumed.
Adjoint frames are orthogonal for the real invariant trace pairing,
and complex orthogonal for its holomorphic bilinear continuation.

At stage \(j\), evaluate \(C_j,\overline C_j,Z_j,B_j,D_j\)
at the ordinary reference
\[
                         U_j=M_\ell^j(U).
\]
Set
\begin{equation}
 \begin{split}
 C^{(t)}&=C_{t-1}\cdots C_0,\\
 \overline C^{(t)}&=\overline C_{t-1}\cdots\overline C_0,\\
 B^{(t)}&=B_{t-1}\cdots B_0,
 \end{split}
 \label{ccs-products}
\end{equation}
with identity products at \(t=0\).
The first is exactly the derivative of the actual ordinary
composition. Define \(Z^{(0)}=0\) and
\begin{equation}
        Z^{(t+1)}=B_tZ^{(t)}+Z_tC^{(t)}.
 \label{ccs-Z-recursion}
\end{equation}

\begin{proposition}[Exact composite tangent identities]
\label{ccs-composite-theorem}
At the final level, writing \(C,\overline C,B,Z,D_c\) for these
products and the final generator,
\begin{equation}
 C-\overline C=D_cZ,\qquad
 \overline C D_{\rm all}=D_cB,\qquad
                  ZD_{\rm all}=r_K-B.
 \label{ccs-composite-identities}
\end{equation}
The operators have exact joint gauge covariance with the original
fine and final ordinary reference frames.
\end{proposition}

\begin{proof}
The primitive intertwining gives
\(\overline C^{(t)}D_0=D_tB^{(t)}\).
Suppose \(C^{(t)}=\overline C^{(t)}+D_tZ^{(t)}\).
Then
\[
 C_tC^{(t)}
  =\overline C_t\overline C^{(t)}
        +D_{t+1}\{B_tZ^{(t)}+Z_tC^{(t)}\},
\]
which proves the first two identities by induction.
For the third, insert
\(Z_tD_t=r_\ell-B_t\) and
\(C^{(t)}D_0=D_tr_{\ell^t}\) in
\eqref{ccs-Z-recursion}; the two intermediate
\(B_tr_{\ell^t}\) terms cancel.
Finally, \eqref{ccs-primitive-covariance} telescopes at the
ordinary reference backgrounds. These are not the
fixed-identity corrected iterates.
\end{proof}

\subsection{Own-block locality and an exact constant row Gram matrix}
\label{ccs-locality-section}

For \(x\in B_K(y)\), let
\(x=x_0,x_1,\ldots,x_k=y\) be its unique nested sequence of
block centers. The scalar product in \eqref{ccs-products}
has only this one ancestry. Consequently
\begin{equation}
 \begin{split}
 B(U;y,x)&=K^{-d}\operatorname{Ad}_{P_{\rm hier}(U;y,x)},\\
 P_{\rm hier}(U;y,x)
  &=P_\ell(U_{k-1};x_k,x_{k-1})\cdots
                                 P_\ell(U_0;x_1,x_0),
 \end{split}
 \label{ccs-hierarchical-B}
\end{equation}
and \(B(U;y,x)=0\) outside \(B_K(y)\).
Each factor is one group element, even when obtained from a
selected group mean of path transports.

Its background dependencies lie in the same final block.
A stage contour uses internal stage bonds. Each such ordinary
bond reads the two descendant blocks at its endpoints, both
inside the parent block by centered nested tiling. Iterating
downwards proves that no coefficient of this scalar row reads
a fine background edge crossing its final block boundary.
The same conclusion holds for all of its derivatives.

The rows of \(C\) and \(\overline C\) read only the two final
endpoint blocks. The scalar \(Z\) row stays in one block:
in \eqref{ccs-Z-recursion}, \(Z_t\) reads internal stage bonds,
and \(C^{(t)}\) on these bonds reads only descendants inside
their parent block. This is full coefficient locality, not
merely support at flatness.

\begin{theorem}[Covariant block projection]
\label{ccs-projector-theorem}
For every real or complex background in the common germ,
\begin{equation}
                 BB^T=K^{-d}I.
 \label{ccs-row-Gram}
\end{equation}
In dimension four,
\begin{equation}
                  P_B=I-K^4B^TB
 \label{ccs-projector}
\end{equation}
is the symmetric projector onto \(\ker B\), orthogonal for
real fields. It is holomorphic, has one-final-block matrix
and background support, and obeys
\begin{equation}
          P_B(U^g)=O_S(g)P_B(U)O_S(g)^T.
 \label{ccs-projector-covariance}
\end{equation}
The normalized term \(\frac12\log\det BB^T\) vanishes identically.
\end{theorem}

\begin{proof}
Distinct rows have disjoint scalar supports because the final
blocks partition the fine sites. In the invariant bilinear
pairing,
\(\operatorname{Ad}_P\operatorname{Ad}_P^T=I\).
Equation \eqref{ccs-hierarchical-B} therefore gives
\[
 (BB^T)_{yy}
      =K^{-2d}\sum_{x\in B_K(y)}I=K^{-d}I,
\]
with zero off-diagonal blocks. This algebra remains valid on
the complexified group. The projector assertions follow
directly, and covariance follows from that of \(B\).
\end{proof}

\subsection{Flat values and the analytic domain}
\label{ccs-analytic-section}

At identity, the contour derivative in \(Z\) is the root term
of the actual ordinary arithmetic differential. Subtraction
of the two endpoint terms in
\eqref{ccs-primitive-correction} leaves precisely
\begin{equation}
 (\overline C_1x)(c)=\ell^{-d}
       \sum_{z\in B_\ell(c_-)}x([z,z+\ell e_\mu]),
 \qquad B_1=Q_{0,\ell}.
 \label{ccs-flat-primitive}
\end{equation}
Centered-block tiling then gives
\begin{equation}
             \overline C(1)=A_K^{\rm st},\qquad
                         B(1)=Q_{0,K}.
 \label{ccs-flat-composite}
\end{equation}
The bond map here is an angle map, with a primitive factor
\(\ell\) relative to the normalized physical-potential map.
There is no such factor for the scalar average.

At fixed \(K\), each coefficient is a composition of local
analytic maps reading finitely many links in one or two final
blocks. A small link-angle product ball therefore works uniformly
in volume and in unrelated preceding depth \(M\).
Its radius and constants may depend on \(K\).
Cauchy derivatives preserve the same support. Compatible tori
are read with the prescribed local lifts and nested partition.

The previously proved ordinary-reference trajectory gives a
stronger optional bound. For \(U=e^{iA/K}\) in its small globally
bounded physical-potential ball, stage angles are
\(O(\ell^j\|A\|_\infty/K)\). Fixed-\(\ell\) local analytic
Schur bounds and the sum
\(\sum_{j<k}\ell^j/K\le(\ell-1)^{-1}\) imply
\begin{equation}
 \|B-Q_{0,K}\|_{{\rm count},\,2\to2}
                    \le C\|A\|_\infty K^{-2}
                    \qquad(d=4).
 \label{ccs-physical-smallness}
\end{equation}
The corresponding counting scale for the angle bond map
\(\overline C\) is \(K^{-1}\).
These stronger estimates are not needed for fixed-\(K\)
volume uniformity or the exact identity
\eqref{ccs-row-Gram}.

The induced scalar map of the fixed-identity corrected
composition is a different operator. It contains logarithm
Jacobians, and its row Gram matrix need not be constant.
The present moving-reference construction avoids that issue;
it must not be replaced by the fixed-identity derivative
without the appropriate endpoint conversion.

\subsection{A basis-free scalar completion}
\label{ccs-completion-section}

Let \(G_{\rm all}=D_{\rm all}^TD_{\rm all}\), and let \(V\)
be any bilinear orthonormal frame of \(\ker B\). Define
\begin{equation}
 G_B=V^TG_{\rm all}V,\qquad
 \widehat G_B=P_BG_{\rm all}P_B+K^{-2}(I-P_B).
 \label{ccs-completion}
\end{equation}
The two projector subspaces reduce \(\widehat G_B\); hence
\begin{equation}
       \det\widehat G_B=\det G_B\,
                        (K^{-2})^{\operatorname{rank}B}.
 \label{ccs-completion-determinant}
\end{equation}
The extra factor is background independent.
The ambient matrix \(\widehat G_B\) is exactly gauge covariant,
without a need to choose gauge-covariant frames \(V\).

For the companion vector matrix
\[
 Q_B=\mathcal H+D_{\rm all}P_BD_{\rm all}^T
                         +a\overline C^T\overline C,
\]
the actual residual covariance and
\eqref{ccs-projector-covariance} give the same conclusion.
When the general tangent determinant identity and the needed
spectral gaps are supplied separately, these observations
permit a formulation using the two ambient matrices
\(\widehat G_B,Q_B\), with no isolated global base site and
no normalized row-Gram term.

Locality and covariance alone do not prove those spectral gaps
or the determinant identity. Once they hold, a fixed-\(K\)
uniform complex perturbation neighborhood permits the usual
diagonal logarithm argument. This does not compute a
large-\(K\) vector Maxwell coefficient or establish a
continuum Yang--Mills theory.


\section{A local block-mean gauge for the actual determinant}
\label{pql-section}

This section first proves a finite-dimensional identity, then continues
it jointly through the flat field. Its scalar gauge has zero block mean,
rather than a point fixed in each block. The resulting scalar and vector
operators have flat gaps of order $K^{-2}$ in current counting units.
The local analytic continuation is uniform in volume and old depth for
each fixed $K$; no uniform large-$K$ curvature comparison is asserted.

\subsection{Actual residual data and an auxiliary coupled matrix}
\label{pql-actual-data}

Fix $K=\ell^k$, old depth $M$, and a compatible finite torus in four
dimensions. Let $E,E_c$ be current and final bond-angle spaces. In this
section $S,S_c$ contain \emph{all} fine and coarse scalar sites, with no
global base removed. Use counting measures and trace-orthonormal color
coordinates. Initially work on a real open set where the left-relative
gauge generators $D:S\to E$ and $D_c:S_c\to E_c$ are injective.
The actual ordinary constraint differential $C$ and center restriction
$r:S\to S_c$ satisfy
\begin{equation}
 CD=D_cr,\qquad D_b=D|_{\ker r}.
 \label{pql-ordinary-intertwining}
\end{equation}
Write the full Lagrangian, coarse and residual Hessians as
\begin{equation}
 L=s''-\lambda\cdot f'',\qquad q=s_{\rm eff}'',\qquad
 \mathcal H=L-C^TqC.
 \label{pql-residual}
\end{equation}
The actual minimum comparison supplies
\begin{equation}
 \mathcal H\ge0,\quad \mathcal H D=0,\quad CH=I,\quad
 \mathcal H H=0.
 \label{pql-minimum-identities}
\end{equation}
On $\ker C$ the nullspace is $\operatorname{ran}D_b$.
The actual contour differential satisfies $TD_b=-I$.
All of these statements retain the constraint-current curvature
$\lambda\cdot f''$.

Let $G=D^TD$, $G_c=D_c^TD_c$, and choose orthonormal bond frames
$U,U_c$ perpendicular to the corresponding gauge ranges. Put
\begin{equation}
 A=U_c^TCU,\quad \mathcal H_\perp=U^T\mathcal H U,\quad
 Q_\perp(a)=\mathcal H_\perp+aA^TA,\qquad a>0.
 \label{pql-physical-auxiliary}
\end{equation}
The actual contour/pivot Gaussian amplitude $I_{\rm act}$ obeys
\begin{equation}
 I_{\rm act}^2
 =\frac{\det G}{\det G_c}
   \frac{a^{n_c-s_c}}{\det Q_\perp(a)},\qquad
 n_c=\dim E_c,\quad s_c=\dim S_c.
 \label{pql-all-site-amplitude}
\end{equation}
Indeed the constrained Gaussian change of variables
$x=Uu+D(\lambda_b,\lambda_c)$ has Jacobian $\sqrt{\det G}$,
while the target split has Jacobian $\sqrt{\det G_c}$.
The target delta fixes $\lambda_c$, and $TD_b=-I$ integrates
$\lambda_b$ with no background determinant. The remaining physical
constraint has $AQ_\perp^{-1}A^T=a^{-1}I$, by
\eqref{pql-minimum-identities}. This proves
\eqref{pql-all-site-amplitude} with the actual target delta in
left-relative coordinates; no additional target Haar factor is inserted.

A preliminary projection-free form is
\begin{equation}
 \mathbb B_a(C)=
 \begin{pmatrix}
 \mathcal H+DD^T+aC^TC&\sqrt a\,C^TD_c\\
 \sqrt a\,D_c^TC&G_c
 \end{pmatrix},\qquad
 I_{\rm act}^2=a^{n_c-s_c}
                   \frac{(\det G)^2}{\det\mathbb B_a(C)}.
 \label{pql-coupled-matrix}
\end{equation}
The ratio is first defined on the injective open set. Its normalized
logarithm is the joint removable amplitude
$\mathcal D=\log I_{\rm act}(W)-\log I_{\rm act}(1)$; the individual
massless logarithms are not evaluated at the flat field. The gapped
formula below supplies their joint continuation. We use
$\mathcal N_0F=F(W)-F(1)$ only for functions defined there.
Its quadratic form is
$x^T\mathcal Hx+|D^Tx|^2+|\sqrt a Cx+D_c\psi|^2$.
Taking the Schur complement in $G_c$ proves
\eqref{pql-coupled-matrix}. Testing $x=0$ shows that this coupled
matrix alone need not have a volume-uniform gap.
If $C=\bar C+D_cZ$, the shear
$\psi'=\psi+\sqrt a Zx$ has determinant one, and hence
\begin{equation}
 \det\mathbb B_a(C)=\det\mathbb B_a(\bar C).
 \label{pql-shear}
\end{equation}
No inverse of $D_c$ is needed for this equality.

\subsection{The exact block-mean determinant}
\label{pql-mean-identity-section}

Assume that $B:S\to S_c$ is onto and that tangent maps satisfy
\begin{equation}
 C-\bar C=D_cZ,\qquad \bar CD=D_cB.
 \label{pql-tangent-identities}
\end{equation}
Neither $\bar C$ nor $B$ is required to be the differential of a
nonlinear pushforward with a prescribed Haar density.
Let $V$ have orthonormal columns spanning $\ker B$, and define
\begin{equation}
 P=VV^T=I-B^T(BB^T)^{-1}B,\quad
 G_B=V^TD^TDV,\quad
 Q_B=\mathcal H+DPD^T+a\bar C^T\bar C.
 \label{pql-mean-operators}
\end{equation}
Then
\begin{equation}
 I_{\rm act}
   =a^{n_c/2}\det G_B\,\det(BB^T)^{1/2}\det Q_B^{-1/2},
 \label{pql-mean-amplitude}
\end{equation}
and therefore
\begin{equation}
 \mathcal D=\mathcal N_0\left[
   \log\det G_B-\tfrac12\log\det Q_B
                          +\tfrac12\log\det(BB^T)\right].
 \label{pql-general-log}
\end{equation}

Here is a direct determinant proof. In orthonormal fine gauge
coordinates $DG^{-1/2}$, the inserted map $DV$ has matrix $G^{1/2}V$,
and the target gauge row of $\bar C$ is $G_c^{1/2}BG^{-1/2}$.
Their product is zero. Split fine gauge coordinates into the range
of $G^{1/2}V$ and its orthogonal complement. The first block contributes
$\det G_B$; elimination of the second block removes the horizontal
coarse-gauge row. Since $P_c\bar C=P_cC$, the remaining physical block
is exactly $Q_\perp(a)$. Thus
\begin{equation}
 \det Q_B=a^{s_c}\det G_B\det G_c
                    \det(BG^{-1}B^T)\det Q_\perp(a).
 \label{pql-vertical-schur}
\end{equation}
Complete $V$ to an orthogonal frame by
$W=B^T(BB^T)^{-1/2}$. The lower inverse block in this frame gives
\begin{equation}
 \det(BG^{-1}B^T)
       =\det(BB^T)\frac{\det(V^TGV)}{\det G}.
 \label{pql-jacobi}
\end{equation}
Combining \eqref{pql-all-site-amplitude},
\eqref{pql-vertical-schur}, and \eqref{pql-jacobi} proves
\eqref{pql-mean-amplitude}. In particular the scalar determinant
has power one, not one half.

There is no need to select a frame of $\ker B$ in the local formula.
Set
\begin{equation}
 \widehat G=P D^TD P+K^{-2}(I-P).
 \label{pql-full-scalar}
\end{equation}
The two $P$ subspaces reduce this operator and
\begin{equation}
 \det\widehat G=(K^{-2})^{s_c}\det G_B.
 \label{pql-scalar-completion}
\end{equation}
Thus $G_B$ may be replaced by $\widehat G$ in
\eqref{pql-general-log}; the extra factor is background-independent.

\subsection{Actual moving-reference tangent maps}
\label{pql-source-maps}

The following construction uses the selected internal contour mean
$P_\ell(U;y,x)$ appearing in the actual ordinary primitive average.
Its needed properties are local analyticity, endpoint covariance,
and arithmetic differential at the identity. No identification with
a different axial contour convention is made.
Use the gauge sign
\[
 (D_U\lambda)(e)=\lambda(e_-)
                  -\operatorname{Ad}_{U(e)}\lambda(e_+).
\]
Let $r_x(U;v)$ be the left differential of $P_\ell(U;y,x)$ in the
bond direction $v$, and define
\begin{equation}
 Z_Uv(y)=\ell^{-d}\sum_{x\in B_\ell(y)}r_x(U;v),\qquad
 B_U\lambda(y)=\ell^{-d}\sum_{x\in B_\ell(y)}
               \operatorname{Ad}_{P_\ell(U;y,x)}\lambda(x).
 \label{pql-primitive-scalar}
\end{equation}
Endpoint covariance gives
$Z_UD_U=r-B_U$. If $C_{{\rm ord},U}$ is the actual ordinary
angle differential and
$C_{{\rm mov},U}=C_{{\rm ord},U}-D_{M_\ell(U)}Z_U$, then
\begin{equation}
 C_{{\rm ord},U}=C_{{\rm mov},U}+D_{M_\ell(U)}Z_U,
 \qquad C_{{\rm mov},U}D_U=D_{M_\ell(U)}B_U.
 \label{pql-primitive-tangent}
\end{equation}
These formulas follow directly from the actual map and its selected
contour mean; they do not modify its nonlinear constraint.

Evaluate these maps at the ordinary intermediate backgrounds
$U_j=M_\ell^j(U)$ and put
\begin{equation}
 \bar C=C_{{\rm mov},k-1}\cdots C_{{\rm mov},0},\qquad
 B=B_{k-1}\cdots B_0.
 \label{pql-composite-tangent}
\end{equation}
Induction in \eqref{pql-primitive-tangent} proves
\eqref{pql-tangent-identities}; the endpoint remainder obeys
$Z_{{\rm tot},j+1}=B_jZ_{{\rm tot},j}+Z_jC_{{\rm ord},j-1}\cdots
C_{{\rm ord},0}$.
Nested blocks give each fine scalar a unique parent chain. Hence
\begin{equation}
 B(y,x)=K^{-d}\operatorname{Ad}_{P_{\rm hier}(U;y,x)}
       \quad(x\in B_K(y)),\qquad BB^T=K^{-d}I.
 \label{pql-exact-scalar-normalization}
\end{equation}
The row supports are disjoint. Adjoint transports preserve the color
inner product, and their bilinear complex continuations preserve the
same identity. At flatness,
\begin{equation}
 B(0)=Q_{0,K},\qquad \bar C(0)=A_K^{\rm st},\qquad
 P(0)=P_{\ker Q_{0,K}}.
 \label{pql-flat-maps}
\end{equation}
The second identity follows by subtracting the two endpoint contour
terms from the ordinary flat differential at each primitive step;
the remaining straight-string averages compose exactly.

Every map above has fixed-$K$ local coefficients. They transform
covariantly under fine and coarse site gauges, because each defining
contour transforms at its two endpoints. Consequently $P$ is a
covariant scalar projection, $\widehat G$ a covariant scalar operator,
and $Q_B$ a covariant bond operator. The scalar normalization in
\eqref{pql-exact-scalar-normalization} is constant. Formula
\eqref{pql-general-log} becomes
\begin{equation}
 \boxed{\displaystyle
 \mathcal D=\mathcal N_0\left(
           \log\det\widehat G-\tfrac12\log\det Q_B\right).}
 \label{pql-covariant-pair}
\end{equation}
This equality still refers to the original ordinary constraint and
its original contour/pivot normalization.

\subsection{Continuation across stabilizers and flat coercivity}
\label{pql-continuation}

For $SU(N)$ the all-site injective locus is nonempty arbitrarily near
flatness. Choose two coarse cycle holonomies $e^{i\epsilon X}$ and
$e^{i\epsilon Y}$, with $X$ distinct diagonal and with the nonzero
offdiagonal graph of the Hermitian matrix $Y$ connected. For small
nonzero $\epsilon$ their common Lie centralizer is zero. Distributing
these holonomies over their cycles gives arbitrarily small coarse
backgrounds with injective $D_c$. If $D\lambda=0$, then
\eqref{pql-ordinary-intertwining} implies $D_cr\lambda=0$;
thus $r\lambda=0$, and parallel transport from one center proves
$\lambda=0$. Injectivity is open at each finite regulator. The same
argument works for compact semisimple groups by choosing generators
with zero common Lie centralizer. An abelian factor requires a separate
stabilizer treatment and is not included in this argument.

At flatness block Poincare gives
\begin{equation}
 G_B(0)\ge cK^{-2}I,\qquad
 \widehat G(0)\ge\min(c,1)K^{-2}I.
 \label{pql-scalar-gap}
\end{equation}
Let $R_q$ be the uniform flat coarse Hessian bound and choose
$a\ge R_q+1$. Since $qD_c=0$ at flatness,
$C^TqC=\bar C^Tq\bar C$ there. The old-depth action dominates
Maxwell, and therefore $Q_B(0)$ dominates
\begin{equation}
 d^*d+D P(0)D^*+(A_K^{\rm st})^*A_K^{\rm st}.
 \label{pql-dominated-form}
\end{equation}
This form has a gap $cK^{-2}$, uniform in $M,K$ and volume.
For a concise proof, use final physical units and decompose
$v=t+D_\eta\phi$ with $D_\eta^*t=0$. Transverse alias Poincare
controls $\norm{t}^2$ by $\norm{d_\eta v}^2+\norm{Qv}^2$.
For $b=Q_0\phi$, the scalar energy-minimizing lift $H_sb$ has
\[
 P\Delta_\eta H_sb=0,\qquad
 \norm{D_\eta H_sb}^2\le C\norm{d_cb}^2.
\]
The latter follows by retaining the central alias in
$(\sum_r|a_r(p)|^2/P_r(p))^{-1}$; its symbol is at most
$P_0/|a_0|^2\le C|d_c(p)|^2$. The zero mode lifts constantly.
Block Poincare applied to $\phi-H_sb$ then gives
\[
 \norm{D_\eta\phi}^2
 \le C\norm{d_cb}^2+C\norm{P\Delta_\eta\phi}^2.
\]
Finally $d_cb=Qv-Qt$ and $Q$ is a physical $L^2$ contraction.
Combining these estimates and restoring counting units proves
\begin{equation}
 Q_B(0)\ge c_aK^{-2}I.
 \label{pql-vector-gap}
\end{equation}
This proves the mean-zero gauge estimate directly; the projection
here is not the earlier source projection onto $\Delta\ker Q_0$.

The operators in \eqref{pql-covariant-pair} are invertible at
flatness, although $D^TD$ and $D_c^TD_c$ are singular there.
Their determinant branches are fixed by positivity. Equality on the
nonempty injective real open set therefore extends holomorphically
through flatness. This is a joint continuation of the finite actual
amplitude, not a limit of its separate massless orbit factors.

\subsection{The fixed-ratio local limit and its remaining scope}
\label{pql-fixed-ratio}

For fixed $K$, the moving-reference coefficients have finite range in the current background $U$. The identity and residual positivity are evaluated on the actual stationary $W$-germ, or its gauge orbit. After substitution $U=U_*(W)$, the source dependence in $W$ is generally exponential, not finite range.
The actual classical support bounds give exponentially local
$\mathcal H$ and its source derivatives, uniformly in old depth and
volume. By \eqref{pql-exact-scalar-normalization},
$P=I-K^dB^TB$ is itself local. The flat gaps permit a common
fixed-$K$ product radius and a positive spatial exponent on which
$\widehat G,Q_B$ have uniformly accretive real parts and bounded
weighted resolvents. Their logarithms and local source derivatives
then have exponentially summable kernels. This follows, for example,
from a resolvent contour around the common positive spectrum and
weighted conjugation; it does not assume a weighted Neumann series
converges merely because an unweighted one does.

Along any fixed jointly locally convergent classical structural
system, compact-source determinant derivatives now pass jointly to
the infinite-volume limit. The finite identity
\eqref{pql-covariant-pair} identifies their total kernels with those
of the actual leading correction before taking that limit. Vacuum
constants are removed by $\mathcal N_0$. No uniqueness of the
classical structural system is inferred. Covariant diagonal-density
and affine-Maxwell consequences can be developed from this gapped
pair separately, with their explicit source and limit conventions.

The radius and weighted source constants in this paragraph may depend
on $K$. Neither the flat gaps nor the exact identity yet prove a
large-$K$ uniform curvature estimate, a bounded reference-comparison
finite part, or the sign of the actual running coupling. The residual
constraint-current term remains present unless a separate on-shell
probe has been established.




% No preamble. Prefix mgs-. Flat mean-zero-gauge spectral estimates.
\section{Both flat mean-zero-gauge factors have a physical-scale gap}
\label{mgs-section}

This section proves spectral estimates for the explicit flat matrices
of the corrected all-site gauge construction.
The exact nonlinear determinant identity and its continuation through
the reducible flat background are separate inputs.
No derivative estimate for a varying-background heat trace is inferred
from these flat bounds.

\subsection{Spaces and flat operators}
\label{mgs-spaces-section}

Work in dimension \(d=4\), on a cubic current torus of side \(KN\),
with \(K=\ell^k\) odd, final centers \(Ky\), and the disjoint centered
cubes \(B_y\) of side \(K\).
All scalar and bond norms are counting norms.
Let \(D\) be the full-site scalar gradient, \(d\) the curl, and
\(\Delta_s=D^TD\).
Let \(Q_0\) be the normalized scalar block average, and put
\begin{equation}
 P=I-Q_0^T(Q_0Q_0^T)^{-1}Q_0,\qquad \Pi=I-P .
 \label{mgs-projections}
\end{equation}
Thus \(P\) projects onto block-mean-zero scalars, and \(\Pi\)
onto block-constant scalars.
Choose a block-diagonal orthonormal frame \(V\) for \(\ker Q_0\)
and define
\begin{equation}
                       G_B=V^T\Delta_sV.
 \label{mgs-ghost}
\end{equation}
At \(K=1\) this scalar space is empty; its trace is zero and
its determinant is one.

The actual flat corrected composite differential is the straight
angle average \(A=A_K\). It obeys
\begin{equation}
 AD=D_cQ_0,\qquad \|A\|_{2\to2}\le K^{-1}.
 \label{mgs-average}
\end{equation}
Here \(D_c\) is the coarse scalar gradient.
The estimate is for this corrected differential; it is not
asserted for the full ordinary unprojected constraint.

Let \(L=q_M\) and \(q=q_{MK}\) be the actual flat classical
Hessians. The established homogeneous inputs are
\begin{equation}
 L\ge d^Td,\qquad LD=0,\qquad 0\le q\le R_qI,
 \label{mgs-classical-inputs}
\end{equation}
uniformly in old depth \(M\), endpoint \(K\), and volume.
Choose a fixed \(a\ge R_q+1\), and set
\begin{equation}
 Q_B=L-A^TqA+DPD^T+aA^TA.
 \label{mgs-vector}
\end{equation}
Its energy satisfies
\begin{equation}
 \langle x,Q_Bx\rangle
 \ge \|dx\|^2+\|PD^Tx\|^2+\|Ax\|^2.
 \label{mgs-energy-lower}
\end{equation}
Matching these matrices to the nonlinear bridge uses only its
flat identifications \(B(1)=Q_0\) and \(\widehat C(1)=A_K\).

\subsection{Scalar Poincare and a block-mean harmonic lift}
\label{mgs-scalar-section}

If \(Q_0\psi=0\), the Neumann cube Poincare inequality, summed
over the blocks, gives
\begin{equation}
 \|\psi\|^2\le C_PK^2\|D\psi\|^2,\qquad C_P=1.
 \label{mgs-poincare}
\end{equation}
Indeed the first nonzero cube eigenvalue is
\(4\sin^2(\pi/(2K))\ge4/K^2\), and cube-internal edges
are a subset of all edges.
The case \(K=1\) is vacuous.
In particular,
\begin{equation}
                         G_B\ge C_P^{-1}K^{-2}I.
 \label{mgs-ghost-gap}
\end{equation}

For a coarse scalar \(b\), let \(H_b\) minimize
\(\|D\phi\|^2\) subject to \(Q_0\phi=b\).
The constraint is onto and fixes the constant nullspace, so the
finite minimizer exists uniquely. Its Euler equation is
\(P\Delta_sH_b=0\).
Moreover,
\begin{equation}
 \|DH_b\|^2\le C_{\rm lift}K^2\|D_cb\|^2,\qquad
 C_{\rm lift}=\frac{\pi^2}{4}\left(\frac\pi2\right)^8.
 \label{mgs-harmonic-upper}
\end{equation}
To verify the normalization, lift a coarse frequency
\(p\in[-\pi,\pi]^4\) to the central fine frequency \(p/K\).
The scalar average has multiplier, up to a unit phase,
\[
 a_0(p)=\prod_{\mu=1}^4
      \frac{\sin(p_\mu/2)}{K\sin(p_\mu/(2K))},
 \qquad |a_0(p)|\ge(2/\pi)^4.
\]
Divide the lifted amplitude by \(a_0\).
The fine difference obeys
\[
 |e^{ip_\mu/K}-1|
       \le\frac{\pi}{2K}|e^{ip_\mu}-1|.
\]
Counting Parseval multiplies the fine squared amplitude by
\(K^4\), relative to its coarse counterpart, proving
\eqref{mgs-harmonic-upper}.
The zero mode is constant and has zero energy.
All selected frequencies are legitimate on the \(KN\)-periodic
torus. For real data one may take the real part of a complex
feasible lift; this preserves the constraint and does not
increase energy, including at coarse Nyquist frequencies.

For arbitrary \(\phi\), put \(b=Q_0\phi\) and
\(\psi=\phi-H_b\). Then
\[
 Q_0\psi=0,\qquad
 P\Delta_s\psi=P\Delta_s\phi=:g,\qquad
 \|D\psi\|^2=\langle\psi,g\rangle.
\]
Equation \eqref{mgs-poincare} therefore gives
\(\|D\psi\|\le\sqrt{C_P}K\|g\|\), and hence
\begin{equation}
 \|D\phi\|^2
 \le 2C_PK^2\|P\Delta_s\phi\|^2
           +2C_{\rm lift}K^2\|D_cQ_0\phi\|^2.
 \label{mgs-scalar-control}
\end{equation}

\subsection{The full vector gap}
\label{mgs-gap-section}

\begin{theorem}[Uniform flat mean-zero-gauge coercivity]
\label{mgs-gap-theorem}
There is a constant \(C_{\rm gap}\), independent of \(M,K,N\),
such that
\begin{equation}
 \|x\|^2\le C_{\rm gap}K^2
       \{\|dx\|^2+\|PD^Tx\|^2+\|Ax\|^2\},
 \qquad
 Q_B\ge C_{\rm gap}^{-1}K^{-2}I .
 \label{mgs-vector-gap}
\end{equation}
\end{theorem}

\begin{proof}
Use the orthogonal periodic Hodge decomposition
\(x=x_T+D\phi\), where \(D^Tx_T=0\).
Harmonic one-forms are included in \(x_T\).
The previously proved transverse alias estimate is
\begin{equation}
 \|x_T\|^2\le K^2
       \{A_0\|P_cAx_T\|^2+B_0\|dx_T\|^2\},
 \quad
 A_0=2c_\omega^{-2},\quad
 B_0=(1+2c_\omega^{-2})/4,\quad c_\omega=(2/\pi)^5.
 \label{mgs-transverse}
\end{equation}
It is a Maxwell/averaging inequality, includes zero momentum,
and has no old-depth dependence.

By \eqref{mgs-average},
\[
 P_cAx_T=P_cAx,\qquad
 \|D_cQ_0\phi\|^2
 \le2\|Ax\|^2+2K^{-2}\|x_T\|^2.
\]
Also \(P\Delta_s\phi=PD^Tx\) and \(dx_T=dx\).
Insert these identities in \eqref{mgs-scalar-control},
then add \(\|x_T\|^2\) and apply \eqref{mgs-transverse}.
A valid constant is
\begin{equation}
 C_{\rm gap}=
 \max\{(1+4C_{\rm lift})B_0,\ 2C_P,\
               (1+4C_{\rm lift})A_0+4C_{\rm lift}\}.
 \label{mgs-gap-constant}
\end{equation}
This proves the first inequality in \eqref{mgs-vector-gap};
\eqref{mgs-energy-lower} proves the second.
\end{proof}

No point-pinned scalar gap or inverse of a globally based scalar
matrix occurs in this argument.

\subsection{A local tent lift controls the full bond Laplacian}
\label{mgs-tent-section}

\begin{lemma}[Block-constant scalar lifting]
\label{mgs-tent-lemma}
There is a block-local map \(E:\operatorname{ran}\Pi\to
\{\text{all scalar sites}\}\) such that
\begin{equation}
 \Pi E=I,\qquad \|E\|\le C_0,\qquad
 \|DE\|\le C_1/K,\qquad
 C_0=2^{d/2},\quad C_1^2=80d\,2^{d-1}.
 \label{mgs-tent-bounds}
\end{equation}
These bounds include all periodic multiplicities.
\end{lemma}

\begin{proof}
For \(K=2h+1\ge3\), put
\[
 f(r)=\frac K{h+1}\left(1-\frac{|r|}{h+1}\right),
 \qquad -h\le r\le h.
\]
Its average is one, its mean square is at most two, and its
endpoint value and adjacent differences are at most \(4/K\).
If \(e\) equals \(c_y\) on \(B_y\), define
\[
 (Ee)(Ky+r)=c_y\prod_{\mu=1}^d f(r_\mu).
\]
The block mean is \(c_y\), and tensorization gives the stated
norm bound. For a fixed direction, internal edges contribute
at most
\(16\,2^{d-1}K^{d-2}\sum_y|c_y|^2\)
to the squared gradient norm.
Interfaces contribute at most
\(64\,2^{d-1}K^{d-3}\sum_y|c_y|^2\), because
\(\sum_y|c_{y+e_\mu}-c_y|^2\le4\sum_y|c_y|^2\).
Divide by
\(\|e\|^2=K^d\sum_y|c_y|^2\) and sum directions.
The same estimates count the repeated periodic interfaces when
\(N=1\) or \(2\).
For \(K=1\), take \(E=I\); the conservative constants still apply.
\end{proof}

For a bond field \(x\) and \(e\in\operatorname{ran}\Pi\),
the equality \(Ee-e\in\operatorname{ran}P\) gives
\[
 \langle\Pi D^Tx,e\rangle
  =\langle x,DEe\rangle-\langle PD^Tx,Ee-e\rangle.
\]
Thus
\begin{equation}
 \|\Pi D^Tx\|
 \le\frac{C_1}{K}\|x\|+(C_0+1)\|PD^Tx\|.
 \label{mgs-duality}
\end{equation}
This step uses only a local lifting and duality.

The full bond Laplacian satisfies
\(\langle x,\Delta_1x\rangle=\|dx\|^2+\|D^Tx\|^2\).
Combining \eqref{mgs-duality},
\eqref{mgs-energy-lower}, and \eqref{mgs-vector-gap}
proves
\begin{equation}
 \begin{split}
 \Delta_1&\le C_\Delta Q_B,\qquad
 C_\Delta=1+2(C_0+1)^2+2C_1^2C_{\rm gap},\\
 Q_B&\ge \frac1{2C_\Delta}\Delta_1
                         +\frac1{2C_{\rm gap}K^2}I.
 \end{split}
 \label{mgs-laplacian-domination}
\end{equation}

\subsection{Uniform heat density and inverse moments for both factors}
\label{mgs-heat-section}

\begin{corollary}[Flat local-gauge heat bounds]
\label{mgs-heat-corollary}
For either \(O=Q_B\) or \(O=G_B\),
\begin{equation}
 N^{-4}\operatorname{Tr}e^{-uK^2O}
 \le C\left\{\frac{K^4}{(1+K^2u)^2}+N^{-4}\right\}e^{-cu},
 \qquad u>0,
 \label{mgs-heat}
\end{equation}
uniformly in \(M,K,N\).
The trace for \(G_B\) is in its block-local
\(\ker Q_0\) frame.
In the infinite flat structural system the \(N^{-4}\) term is
omitted. Both factors obey the inverse-moment bounds and the
weighted final-time exponential tail of
Theorem~\ref{pvs-main}, with \(O\) in place of its physical \(Q\).
\end{corollary}

\begin{proof}
For \(Q_B\), apply eigenvalue min--max to
\eqref{mgs-laplacian-domination}.
The comparison is the full scalar lattice heat trace, with a
fixed rescaling of time and multiplicity \(4n_{\mathfrak g}\),
times \(e^{-cs/K^2}\).

For \(G_B\), \eqref{mgs-ghost-gap} gives
\[
 G_B\ge\tfrac12G_B+(2C_PK^2)^{-1}I.
\]
The eigenvalues of the compression \(V^T\Delta_sV\)
dominate the correspondingly indexed full scalar eigenvalues,
by the variational principle.
Its heat trace is therefore at most the full scalar heat trace.
The gap supplies the same exponential factor.

The scalar heat estimate and the substitution \(s=K^2u\) are
exactly those used in Theorem~\ref{pvs-main}.
Final-block translation invariance identifies the normalized
total with one block trace.
Coarse Bloch integration proves the infinite version directly.
Mellin integration gives the stated inverse powers and tails.
\end{proof}

The basis-free scalar completion used by the analytic determinant
bridge is
\begin{equation}
 \widehat G_B=P\Delta_sP+K^{-2}(I-P).
 \label{mgs-basis-free-scalar}
\end{equation}
At flatness its two subspaces reduce it. Its heat density in final
time is exactly
\[
 N^{-4}\operatorname{Tr}e^{-uK^2\widehat G_B}
 =N^{-4}\operatorname{Tr}e^{-uK^2G_B}
                        +n_{\mathfrak g}e^{-u}.
\]
This extra term satisfies \eqref{mgs-heat}, after decreasing \(c\):
for \(u\le1\),
\(K^4/(1+K^2u)^2\ge1/4\), while for \(u\ge1\) it is at least
\(1/(4u^2)\), and \(u^2e^{-(1-c)u}\) is bounded when \(c<1\).
The normalized inverse moment adds exactly \(n_{\mathfrak g}\),
which is absorbed in every case of \eqref{pvs-moments}.
Thus the actual basis-free pair \(\widehat G_B,Q_B\) has the
same heat, inverse-moment, and final-time tail bounds.

These estimates concern the explicit flat local-gauge factors.
They do not establish uniform relative-form derivatives under an
affine curvature source, or identify separate massless factors
with independent continuum references.
The exact nonlinear determinant bridge and the remaining
large-\(K\) curvature comparison must retain their own
hypotheses and proofs.



% No preamble. Prefix cgd-. Depends on the checked covariant tangent bridge.
\section{Actual local densities in the covariant mean-zero gauge}
\label{cgd-section}

Use the covariant tangent construction above, with all-site
fine gauge generator \(D\), covariant scalar block average
\(B\), corrected bond differential \(\bar C\), and
\[
 BB^T=K^{-4}I,\qquad
 P=I-K^4B^TB,\qquad
 \widehat G=P D^TD P+K^{-2}(I-P),\qquad
 Q_B=\mathcal H+D P D^T+a\bar C^T\bar C.
\]
The actual finite determinant identity is
\begin{equation}
 D_{M,K}=\mathcal N_0
       \left(\log\det\widehat G-\tfrac12\log\det Q_B\right).
 \label{cgd-identity}
\end{equation}
Here \(\mathcal H\) is the full residual Hessian, not a
substitution of the bare Wilson Hessian for \(S_M''\).
The matrix \(P\) projects onto scalar fields with zero
covariant block mean. It is not the projection onto
fields zero at the block centers.

\subsection{Gauge covariance and localized logarithms}
\label{cgd-covariance-section}

Let a fine site gauge induce orthogonal adjoint matrices
\(O_S,O_E,O_c\) on fine scalars, fine bonds, and coarse
scalars. The covariant tangent construction gives
\[
 B(U^g)=O_c B(U)O_S^T,\qquad
 P(U^g)=O_S P(U)O_S^T,\qquad
 \bar C(U^g)=O_{E_c}\bar C(U)O_E^T.
\]
Together with the exact covariance of \(D,\mathcal H\)
this proves
\begin{equation}
 \widehat G(U^g)=O_S\widehat G(U)O_S^T,\qquad
 Q_B(U^g)=O_EQ_B(U)O_E^T.
 \label{cgd-covariance}
\end{equation}
These are also complex-orthogonal identities on the
holomorphic germs. No choice of a frame for \(\ker B\)
appears in them.

Fix \(K\). The flat lower bounds for both matrices are
\(cK^{-2}\), uniformly in old depth and volume. Their
flat upper bounds and exponentially localized kernels,
and their first-two source kernels after actual
stationary substitution, are uniform in those same
parameters at this fixed \(K\).
The block-local scalar projection does not enlarge
support beyond the corresponding final block; the
bond construction has bounded range in final units.
The old action and stationary response retain their
established exponential source bounds.

Consequently there is a complex output ball, allowed
to depend on \(K\), on which both logarithms have
source-rooted exponential kernels uniformly in volume
and old depth. One direct proof is the logarithm series.
Choose \(s_0\) larger than the flat upper spectral bound.
At flatness \(\|I-T/s_0\|<1\), for either matrix \(T\).
A sufficiently small exponential conjugation and
background perturbation preserve that strict bound.
Then
\[
 \log T=(\log s_0)I-
       \sum_{n\ge1}\frac{(I-T/s_0)^n}{n}
\]
converges in the resulting weighted matrix algebra.
First and second source derivatives cost at most
polynomial factors in \(n\), absorbed by the geometric
decay. Every source remains attached to the matrix
product chain. Spending a smaller exponential converts
the weighted bounds into absolute rooted sums.
The fine fiber dimension is finite at this fixed \(K\).
No endpoint-independent source constant is inferred
from this argument.

Group scalar sites and positive fine bonds by their
initial final block \(x\), and define
\begin{equation}
 \rho_x(W)=\mathcal N_0\left\{
  \operatorname{Tr}_{x,S}\log\widehat G(W)
       -\tfrac12\operatorname{Tr}_{x,E}\log Q_B(W)\right\}.
 \label{cgd-density}
\end{equation}
The finite sum is \(D_{M,K}\). Each diagonal color
trace is individually gauge invariant by
\eqref{cgd-covariance}. The actual stationary branch
under an output gauge changes by a fine gauge; its
root-gauge convention therefore does not change these
traces. Along the chosen jointly locally convergent structural
system, these densities have local infinite limits
and absolutely summable polynomial source moments,
without introducing a globally based scalar inverse.

In particular, each of the two new factors separately
has zero thermodynamic quadratic response to a
constant commuting connection, at fixed \(K\).
Apply its local Ward identity to a linear scalar
parameter and pass the absolutely summable Hessian
row sum. Thus the growing constant-connection terms
of the globally based physical quotient have
already canceled in the exact change of variables;
they are not discarded by a subsequent limit.

\subsection{Matching the actual Maxwell coefficient}
\label{cgd-Maxwell-section}

\begin{proposition}[Local density and actual affine coefficient]
\label{cgd-density-theorem}
In every joint infinite structural system at fixed \(K\),
the density in \eqref{cgd-density} satisfies
\begin{equation}
 \beta(D_{M,K})=
    2[t^2]\rho_0\bigl(w_\mu(n)=t n_\nu H\bigr),
 \qquad \mu\ne\nu,\quad\tau(H^2)=1.
 \label{cgd-Maxwell}
\end{equation}
The polynomial substitution uses the first two localized
Taylor kernels. It is not an evaluation of a bounded-field
analytic map at a nonzero unbounded connection.
\end{proposition}

\begin{proof}
The linear color jet of each invariant density is zero
for the semisimple color algebra. Its quadratic Ward
identity kills gradients in either slot. Exponential
source moments extend the identity to linear scalar
parameters, so its row and column sums vanish for
each fixed bond component. Write \(k_0(b,c)\) for that
density's Hessian on the chosen component and color.
Then
\[
 [t^2]\rho_0(t n_\nu H)
   =\frac12\sum_{b,c}k_0(b,c)b_\nu c_\nu
   =-\frac14\sum_{b,c}k_0(b,c)(b_\nu-c_\nu)^2.
\]
All sums are absolute. Translation covariance and
reanchoring identify the last expression with
\(-\frac14\sum_r r_\nu^2\mathcal K_{\mu\mu}(0,r)\),
where \(\mathcal K\) is the actual total Hessian
kernel. Its Maxwell moment is
\(-\frac12\sum_r r_\nu^2\mathcal K_{\mu\mu}(0,r)\).
This proves \eqref{cgd-Maxwell}.
\end{proof}

The actual stationary polynomial jets on this probe
are the commuting affine fine connection established
in Notebook 33. All gauge-covariant matrices may
therefore be computed in that representative.
The full effective classical Hessian and the
covariant averaging operators remain in \(Q_B\);
the formula does not replace their curvature
derivatives by a continuum operator.

\subsection{A coupled heat representation with the physical flat scale}
\label{cgd-heat-section}

Let \(\widehat G(t),Q_B(t)\) denote the actual affine
jets just specified, and set
\begin{equation}
 h^{\rm mean}_{M,K}(s)
  =[t^2]\left\{
    \operatorname{Tr}_{0,E}e^{-sQ_B(t)}
             -2\operatorname{Tr}_{0,S}e^{-s\widehat G(t)}
         \right\}.
 \label{cgd-heat-coefficient}
\end{equation}
The fixed-\(K\) logarithm/resolvent identities, or
ordered Duhamel with the localized first-two vertices,
give the exact representation
\begin{equation}
 \boxed{\displaystyle
     \beta(D_{M,K})=
        \int_0^\infty h^{\rm mean}_{M,K}(s)\,\frac{ds}{s}.}
 \label{cgd-heat-identity}
\end{equation}
There is no additional coarse determinant in this
covariant choice: \(BB^T=K^{-4}I\), and the added
scalar complement in \(\widehat G\) contributes only
a background-independent determinant.

For precision, this is a fixed-\(K\) identity on
localized Taylor kernels. The spectral gaps,
exponential weighted kernels and finite polynomial
vertices control the first-two Duhamel coefficients
by an integrable small-time and exponentially
decaying large-time envelope. Its constants may
depend on \(K\). This justifies the logarithm
integral, polynomial source substitution and local
limit; it does not give a uniform curvature
estimate as \(K\) grows.

The new scalar and vector factors both have flat
spectral scale \(K^{-2}\). The point-pinned ghost
coefficient \(4\alpha_g\log K+O(1)\) of Notebook 34
belongs to a different exact determinant splitting
and must not be assigned to \(\widehat G\).
The sum in \eqref{cgd-heat-identity} is the same
actual coefficient. Its endpoint-uniform magnetic
comparison remains to be proved in these new,
local, gauge-covariant variables.


\section{The remaining comparison}

Notebook 34 proves the complete based-ghost Maxwell
contribution \(4\alpha_g\log K+O(1)\) and the full
actual ultraviolet bound. In its original local
gauge choice, the remaining vector coefficient
must be \(18\alpha_g\log K+O(1)\) to obtain the
specified continuum coefficient
\(22\alpha_g\log K+O(1)\).

The physical quotient remains an exact finite
identity. Its separate massless factors cannot be
replaced by independently regulated determinants:
the constant-connection calculation gives an
explicit noncommutation and the compensation
that such a replacement would discard.

The new covariant block-mean representation resolves
that infrared issue algebraically. It is an exact
identity for the actual finite amplitude, followed
by a common local analytic continuation and a
fixed-\(K\) local limit. Both factors have the
physical flat spectral scale, and their diagonal
densities have individual gauge identities.
No separately singular globally based factor
remains in the new heat formula.

The next estimate is the uniform large-\(K\)
magnetic comparison for this actual pair. The
flat heat bounds do not by themselves bound
its curvature insertions. The earlier point-pinned
ghost coefficient belongs to the old splitting;
the new scalar coefficient must be computed
with its covariant mean projection retained.
The required total is still
\(22\alpha_g\log K+O(1)\), uniformly in preceding
depth and the required structural systems.

A positive actual block drift would activate the
finite transfer and tuning results of Notebook 32.
Full Yang--Mills existence also requires the
measure matching, mixed and large-field histories,
sourced continuum limits, nontriviality, Euclidean
covariance, positivity and reconstruction in the
construction notebook. Those requirements remain
part of the same objective.

\begin{thebibliography}{9}
\bibitem{priorir}
\emph{Toward a four-dimensional Yang--Mills construction},
Notebooks 32--34 in this workspace.
\bibitem{sourceir}
E.~Sussman and collaborators,
\href{https://github.com/ethanwsussman/Yang--Mills-LLM-experiment}
{Yang--Mills-LLM-experiment}, deposited accounts pinned at
commit \texttt{e392ef4501\allowbreak 3a581664ba\allowbreak b76b908671\allowbreak 70b30fe10e}.
\end{thebibliography}
\end{document}
